% !PRINTED 169
\weberpage{169}

%<<CONT
are. By IV, $P$ would then have to divide the divisor of $T$, and
$F(t)$ and $f(t)$ could not be relatively prime. By \S.\,6 we can
therefore determine two integral rational functions $Q$, $q$ of $x$
and $t$ and an integral rational function $X$ of the $x$ alone in such a way
that the identity

\weq{4}{QF(t) + qf(t) = X}
holds.

Now if $\Phi(t)$ is a further integral rational function of $x$
and $t$, then from (4), by multiplication with $\Phi(t)$, it follows that
\[
Q\Phi(t)F(t) + q\Phi(t)f(t) = X\Phi(t),
\]
and hence, if $\Phi(t)F(t)$ is divisible by $f(t)$, then by this
$X\Phi(t)$ is also divisible by $f(t)$.

Accordingly, if $\varphi(t)$ and $\psi(t)$ again denote two integral
functions of $x$ and $t$, we may set

\weq{5}{\begin{aligned} X\Phi(t) &= \varphi(t)f(t) \\ F(t)\Phi(t) &= \psi(t)f(t). \end{aligned}}

If one multiplies the second of these equations by $X$ and
substitutes from the first the expression $\varphi(t)f(t)$ for $X\Phi(t)$, then
$f(t)$ can be cancelled and it follows that
\[
\varphi(t)F(t) = X\psi(t).
\]
By IV., $X$ must therefore divide both the divisor of $\varphi(t)f(t)$ and
the divisor of $\varphi(t)F(t)$, and since the functions $f(t)$
and $F(t)$, and consequently also their divisors, are relatively prime, $X$ must
divide the divisor of $\varphi(t)$. If we accordingly set
\[
\varphi(t) = X\varphi_1(t),
\]
then from (5) it follows that
\[
\Phi(t) = \varphi_1(t)f(t),
\]
that is, $\Phi(t)$ is divisible by $f(t)$; hence:

\gesperrt{VII. If $F(t)$ and $f(t)$ are relatively prime and $\Phi(t)F(t)$
is divisible by $f(t)$, then $\Phi(t)$ is divisible by $f(t)$.}

But this is nothing other than Theorem I for functions
of $n + 1$ variables, and I. is thus proved in general.
